\begin{figure}[tbp]
\centering
\begin{tikzpicture}[
  stage/.style={draw=MidnightBlue!65,fill=MidnightBlue!3,
    rounded corners=1pt,align=center,font=\small,
    text width=54mm,minimum height=11mm,inner sep=3pt},
  arrow/.style={-{Latex[length=2mm]},draw=MidnightBlue!80,thick}
]
\node[stage,text width=67mm] (data) at (0,0)
  {Full trace cone\\Realification and ideal transport (Section~\ref{sec:structure})};
\node[stage,text width=67mm] (errors) at (0,-1.6)
  {Fixed-cut comparisons and local error ideals\\(Section~\ref{sec:sequence-models})};
\node[stage] (small) at (-3.5,-3.3)
  {Small stabilizers, with or without\\a given model (Section~\ref{sm:section})};
\node[stage] (rough) at (3.5,-3.3)
  {Tracial approximation on the cone\\(Section~\ref{tc:section})};
\node[stage] (unique) at (-3.5,-5.0)
  {Norm uniqueness for given models\\(Section~\ref{sec:uniqueness})};
\node[stage] (cone) at (3.5,-5.0)
  {Fine-label correction and\\retained lower support (Section~\ref{cc:section})};
\node[stage,text width=67mm] (point) at (0,-6.8)
  {Covariant line model and point model\\(Section~\ref{pt:section})};
\node[stage,text width=67mm] (actual) at (0,-8.4)
  {Actual homomorphisms and intertwining\\(Section~\ref{pt:section})};
\draw[arrow] (data) -- (errors);
\draw[arrow] (errors.south west) -- (small.north);
\draw[arrow] (errors.south east) -- (rough.north);
\draw[arrow] (small) -- (unique);
\draw[arrow] (rough) -- (cone);
\draw[arrow] (small.east) -- (0,-3.3) -- (0,-5.0) -- (cone.west);
\draw[arrow] (unique.south) -- (point.north west);
\draw[arrow] (cone.south) -- (point.north east);
\draw[arrow] (point) -- (actual);
\draw[arrow,dashed] (unique.west) -- (-7.0,-5.0) -- (-7.0,-8.4) -- (actual.west);
\end{tikzpicture}
\caption{Principal model-construction dependencies. The existence and
uniqueness arguments meet in the passage from sequence models to actual
maps. Shared auxiliary inputs are not drawn separately: the trace-cap
and barycenter lemmas in Section~\ref{sec:uniqueness} are used in
Section~\ref{tc:section}; the absorption and hereditary-support lemmas
of Section~\ref{sec:uniqueness} are used in Section~\ref{cc:section}.
Norm uniqueness is a statement about
models when they are given; its proof does not assume the cone-existence
theorem. The dashed arrow records its further use in actual-map
extraction and the final two-sided intertwining.}
\label{fig:proof-roadmap}
\end{figure}
